\section{Introduction} \label{sec:intro}

Deep generative models have significantly improved in the past few years \cite{biggan, gpt2, oord2016wavenet}. This is, in part, thanks to architectural innovations as well as computation advances that allows training them at larger scale in both amount of data and model size. The samples generated from these models are hard to distinguish from real data without close inspection, and their applications range from super resolution \cite{srgan} to domain editing \cite{cycleGan}, artistic manipulation \cite{dtn}, or text-to-speech  and music generation \cite{oord2016wavenet}. 

\begin{figure}[hb!]
\centering
\includegraphics[width=0.32\textwidth]{samples_jpg/im256_3.jpeg}
\includegraphics[width=0.32\textwidth]{samples_jpg/im256_2.jpeg}
\includegraphics[width=0.32\textwidth]{samples_jpg/im256_1.jpeg}
\caption{Class-conditional 256x256 image samples from a two-level model trained on ImageNet.} 
\label{fig:samples}
\end{figure}

We distinguish two main types of generative models: likelihood based models, which include VAEs \cite{VAE, Rezende2014}, flow based \cite{dinh2014nice, rezende2015variational, dinh2016density, kingma2018glow} and autoregressive models \cite{larochelle2011neural, PixelRNN}; and implicit generative models such as Generative Adversarial Networks (GANs) \cite{goodfellow2014generative}. Each of these models offer several trade-offs such as sample quality, diversity, speed, etc. 

GANs optimize a minimax objective with a generator neural network producing images by mapping random noise onto an image, and a discriminator defining the generators' loss function by classifying its samples as real or fake. Larger scale GAN models can now generate high-quality and high-resolution images \cite{biggan, karras2018style}. However, it is well known that samples from these models do not fully capture the diversity of the true distribution. Furthermore, GANs are challenging to evaluate, and a satisfactory generalization measure on a test set to assess overfitting does not yet exist. For model comparison and selection, researchers have used image samples or proxy measures of image quality such as Inception Score (IS) \cite{salimans2016improved} and Fr\'echet Inception Distance (FID) \cite{FID}.

In contrast, likelihood based methods optimize negative log-likelihood (NLL) of the training data. This objective allows model-comparison and measuring generalization to unseen data. Additionally, since the probability that the model assigns to \emph{all} examples in the training set is maximized, likelihood based models, in principle, cover all modes of the data, and do not suffer from the problems of mode collapse and lack of diversity seen in GANs. In spite of these advantages, directly maximizing likelihood in the pixel space can be challenging. First, NLL in pixel space is not always a good measure of sample quality \cite{Theis2016a}, and cannot be reliably used to make comparisons between different model classes. There is no intrinsic incentive for these models to focus on, for example, global structure. Some of these issues are alleviated by introducing inductive biases such as multi-scale \cite{theis2015generative, PixelRNN, reed2017parallel, SPN} or by modeling the dominant bit planes in an image \cite{kolesnikov2017pixelcnn, kingma2018glow}.



In this paper we use ideas from lossy compression to relieve the generative model from modeling negligible information. Indeed, techniques such as JPEG \cite{wallace1992jpeg} have shown that it is often possible to remove more than 80\% of the data without noticeably changing the perceived image quality. As proposed by \cite{VQVAE}, we compress images into a discrete latent space by vector-quantizing intermediate representations of an autoencoder. These representations are over 30x smaller than the original image, but still allow the decoder to reconstruct the images with little distortion. The prior over these discrete representations can be modeled with a state of the art PixelCNN \cite{PixelRNN, pixelcnn} with self-attention \cite{Vaswani2017}, called PixelSnail \cite{PixelSnail}. When sampling from this prior, the decoded images also exhibit the same high quality and coherence of the reconstructions (see \figref{fig:samples}). Furthermore, the training and sampling of this generative model over the discrete latent space is also 30x faster than when directly applied to the pixels, allowing us to train on much higher resolution images. Finally, the encoder and decoder used in this work retains the simplicity and speed of the original VQ-VAE, which means that the proposed method is an attractive solution for situations in which fast, low-overhead encoding and decoding of large images are required.